\section{Background and Related Works}
\label{sec:background}
\subsection{Contextual Integrity}

Contextual Integrity (CI) defines privacy as the appropriateness of an
\textit{information flow} to its \textit{context}~\cite{nissenbaum}. Each flow
is described by five parameters: Data Subject, Sender, Recipient, Information
Type, and Transmission Principle. A flow respects contextual integrity when
\textit{all} five conform to the norms of the prevailing context.

Barth et al.~\cite{barth2006} adopted CI for computer science, expressing norms
derived from legislation in linear temporal logic; the norm is an input to
their framework. Later work has instead elicited norms from people, through
vignette surveys and factorial designs that measure what a population judges
appropriate~\cite{martin2016,apthorpe2018}, including for LLM
chatbots~\cite{tran2025}.

\subsection{Retrieval and Agents}
Retrieval-augmented generation conditions a model on documents fetched from an
external store at inference time~\cite{lewis2020}, and tool-assisted agents
extend this into a loop in which the model chooses when to query, what to query
for, and how to act on the result~\cite{yao2023,schick2023}. The set of
documents an agent will touch is therefore not known in advance. Models can
also leak memorized training data~\cite{carlini2021}, and retrieval corpora add
leakage channels of their own~\cite{zeng2024}.

Ruan et al.~\cite{ruan2024} construct sandboxed emulations to surface the risks
agents realize when acting with real permissions. Others show the instruction
channel itself is unreliable: models can be diverted by adversarial
instructions in their prompts~\cite{perez2022}, including instructions embedded
in the documents they retrieve~\cite{greshake2023}. We read these findings as
an argument about adjudicator placement: each describes a failure to comply
with a constraint the model was asked to honor. Recent prompt-injection
defenses also move enforcement outside the model, applying capability or
information-flow policies to what an agent may do with the data it
reads~\cite{debenedetti2026camel,costa2025fides}.

\subsection{Contextual Integrity for Language Models}
A growing body of work applies contextual integrity to language models
directly. ConfAIde evaluates whether models respect contextual norms when
handling information disclosed in conversation, finding that they frequently
reveal what context forbids~\cite{mireshghallah2024}. PrivacyLens extends this
from judgment to action, measuring norm violations in agent
trajectories~\cite{shao2024}. Later benchmarks extend the evaluation to
assistants on synthetic data~\cite{cheng2024cibench}, to persistent
memory~\cite{mireshghallah2026cimemories}, and to enterprise agents working from
retrieved documents, where larger models and deeper reasoning do not solve the
problem~\cite{fu2026ciwork}. GoldCoin grounds the evaluation in law,
generating cases from HIPAA provisions to test whether models reason correctly
about statutory privacy~\cite{goldcoin}; Privacy Checklist uses language models
to cover HIPAA's regulations in a checklist for detecting
violations~\cite{li2025checklist}; and PrivaCI-Bench assembles flows from real
court decisions so that the label authority is a ruling~\cite{privacibench}.
Model judgments of CI norms change with small variations in the
prompt~\cite{shvartzshnaider2024llmci}; we find they change even when the
prompt is identical (Section~\ref{sec:selfconsistency}). Training models to
reason about CI reduces inappropriate disclosure~\cite{lan2025}, but the model
still makes the decision.

A second line proposes architecture. AirGapAgent restricts the context
available to a conversational agent so that it cannot disclose what it was
never given~\cite{bagdasarian2024}, and Ghalebikesabi et
al.~\cite{ghalebikesabi2024} operationalize contextual integrity in assistants
by having the system reason explicitly about the appropriateness of each
disclosure. These are the closest antecedents to our work, and we share their
premise that limiting what reaches the model is more reliable than constraining
what it does afterward.

They differ from our proposal in what the model is asked to judge. In each case
a model inference determines the normative question of who ought to receive the
document, so the decision inherits the variance of the component making it. We
instead require that the normative judgment be fixed in advance by an
accountable policy authority, in a label the enforcement layer reads
mechanically.

\subsection{Privacy Conductor vs.\ RBAC, ABAC, and Conventional Policy Enforcement}
The mechanism we use is not new. Role-based access control has governed
enterprise data for decades~\cite{sandhu1996}, attribute-based access control
generalizes it to arbitrary properties of the subject, resource, and
environment~\cite{hu2014abac}, and our gate is recognizably a policy decision
point in that tradition. Retrieval pipelines already filter documents by role
before the model sees them~\cite{chen2025sac}, while Permissioned LLMs enforce
access control inside the model, with adapters fine-tuned per access
domain~\cite{jayaraman2025permissioned}.

\begin{table}[!t]
\centering
\caption{Privacy Conductor (PC) compared with conventional policy enforcement
and LLM-mediated Contextual Integrity (LLM-CI).}
\label{tab:comparison}
\footnotesize
\setlength{\tabcolsep}{3pt}
\resizebox{0.82\columnwidth}{!}{\begin{tabularx}{\columnwidth}{@{}>{\raggedright\arraybackslash}Xllll@{}}
\toprule
\textbf{Property} & \textbf{RBAC} & \textbf{ABAC} & \textbf{LLM-CI} & \textbf{PC} \\
\midrule
Role-aware                         & Yes     & Yes            & Potentially & Yes \\
Subject relationship               & Limited & Yes            & Potentially & Yes \\
Contextual transmission principle  & No      & Expressible    & Yes         & Explicit \\
Deterministic decision             & Yes     & Yes            & No          & Yes \\
LLM inside policy decision         & No      & No             & Usually     & No \\
Request text affects authorization & Rarely  & Depends        & Often       & No \\
CI-oriented audit explanation      & No      & Not inherently & Potentially & Yes \\
\bottomrule
\end{tabularx}
}
\end{table}

Table~\ref{tab:comparison} gives the comparison. The gate combines a role
condition with a relationship condition~\cite{fong2011}
(Section~\ref{sec:gate}), so it conforms to its specification by construction.
What differs is placement and provenance: it sits inside the retrieval path,
never reads the request text, and enforces a policy written by an accountable
authority. The novelty is the separation of model reasoning from normative
privacy authorization.
